\chapter{Problem Statement and Scope}

\section{Problem Statement}
Existing computing models manage state persistently across user sessions, creating a multitude of issues in high-security, educational, or highly private situations:
\begin{table}[h]
    \centering
    \begin{tabularx}{\textwidth}{lX}
    \toprule
    \textbf{Limitation} & \textbf{Vanish Lite Resolution} \\
    \midrule
    State Persistence & Implements strict, resilient teardown guarantees enforcing zero local persistence through \texttt{userdel} and process termination hooks. \\
    Resource Heavy Isolation & Operates inherently on the host OS by using Linux user paradigms and local group management instead of running heavy virtual machines. \\
    Static Security Modeling & Provides dynamic mode-based policies (\texttt{dev}, \texttt{secure}, \texttt{privacy}, \texttt{exam}, \texttt{online}) configurable down to DNS resolution and USB access constraints. \\
    Complex Auditing & Implements automated per-session snapshots and compliance reports, seamlessly integrated with a management web interface. \\
    \bottomrule
    \end{tabularx}
    \caption{Limitations of traditional computing paradigms compared to Vanish Lite.}
    \label{tab:limitations}
\end{table}

The primary technological hurdle lies in engineering an extensible session enforcement mechanism that can transparently intercept operations and apply strict quotas or network allowances without requiring deep kernel module dependencies.

\section{Objectives}
The core objectives of the Project Vanish Lite platform are to:
\begin{enumerate}
    \item Design an engine that correctly performs dynamic, error-free provisioning and deprovisioning of temporary Linux user accounts.
    \item Implement tailored modes restricting CPU, memory, Network access, and specific executables using local firewall (\texttt{iptables}) commands and \texttt{cgroups}-like constraints.
    \item Create a failsafe cleanup daemon guaranteeing the absolute removal of all temporary processes and persona configurations, preventing orphan environments.
    \item Develop an intuitive Admin Dashboard providing total visibility, telemetry generation, and snapshot reviews.
\end{enumerate}

\section{Scope}
The scope of the project encompasses the development, deployment, and testing of the Linux session management engine and the web-based administrative dashboard. It natively encompasses the following components:
\begin{itemize}
    \item \textbf{C++ Enforcement Engine:} Handles system calls, handles \texttt{sigaction} loops, and invokes network rule modifiers securely.
    \item \textbf{Mode Policies:} \texttt{secure} (restricting peripheral input), \texttt{exam} (whitelisting network rules, mounting persistence drives), \texttt{privacy} (RAM-mounted home directories), and \texttt{online} (configurable dynamic networking paths).
    \item \textbf{Server API:} A Python-backend orchestrating the user interface, session history reporting, and configurations.
\end{itemize}
The current iteration targets modern Linux distributions, assuming root execution privileges for the central daemon.
